% ===== PRÉAMBULE CANONIQUE — à coller VERBATIM en tête de CHAQUE .tex =====
\documentclass[11pt,a4paper]{article}
\usepackage[utf8]{inputenc}\usepackage[T1]{fontenc}\usepackage{lmodern}
\usepackage[french]{babel}
\usepackage{amsmath,amssymb,mathtools}\usepackage{icomma}
\usepackage{siunitx}\sisetup{locale=FR}  % \qty{}{}, \si{}, \ang{}
\usepackage{xcolor}\usepackage{colortbl}
\usepackage{tikz}\usepackage{pgfplots}\pgfplotsset{compat=1.18}
\usepackage[siunitx,europeanresistors]{circuitikz} % circuits (élec.)
\usepackage[most]{tcolorbox}\usepackage{enumitem}\usepackage{array,booktabs}
\usepackage[a4paper,margin=2.2cm]{geometry}
\usetikzlibrary{arrows.meta,calc,angles,quotes,positioning,patterns,%
  backgrounds,shapes.geometric,decorations.pathmorphing,decorations.markings,intersections}
\definecolor{primary}{RGB}{222,49,99}\definecolor{primarydark}{RGB}{150,22,55}
\definecolor{defcol}{RGB}{0,114,178}\definecolor{methcol}{RGB}{52,92,148}
\definecolor{excol}{RGB}{0,135,90}\definecolor{attncol}{RGB}{200,40,55}
\tcbset{boxbase/.style={enhanced,breakable,boxrule=0.9pt,arc=2.5pt,
  left=3mm,right=3mm,top=2.3mm,bottom=2.3mm,fonttitle=\bfseries,coltitle=white}}
% --- boîtes TUTORIEL (noms EXACTS, ne pas renommer) ---
\newtcolorbox{cle}[1][]{boxbase,colback=defcol!6,colframe=defcol,
  title={\textsc{À retenir}\ifx\relax#1\relax\relax\else\textmd{ : \itshape #1}\fi}}
\newtcolorbox{methode}[1][]{boxbase,colback=methcol!7,colframe=methcol,sharp corners,
  title={\textsc{Méthode}\ifx\relax#1\relax\relax\else\textmd{ --- \itshape #1}\fi}}
\newtcolorbox{exemplebox}[1][]{boxbase,colback=excol!5,colframe=excol,
  title={\textsc{Exemple}\ifx\relax#1\relax\relax\else\textmd{ : \itshape #1}\fi}}
\newtcolorbox{piege}[1][]{boxbase,colback=attncol!6,colframe=attncol,
  title={\textsc{Piège}\ifx\relax#1\relax\relax\else\textmd{ : \itshape #1}\fi}}
\newtcolorbox{remarque}[1][]{boxbase,colback=black!4,colframe=black!45,
  title={\textsc{Remarque}\ifx\relax#1\relax\relax\else\textmd{ : \itshape #1}\fi}}
% --- boîtes EXERCICE / CORRIGÉ (noms EXACTS, auto-comptées) ---
\newtcolorbox[auto counter]{exo}[1][]{boxbase,colback=primary!4,colframe=primary,
  title={\textsc{Exercice}~\thetcbcounter\ifx\relax#1\relax\relax\else\textmd{ --- \itshape #1}\fi}}
\newtcolorbox[auto counter]{corrige}[1][]{boxbase,colback=excol!4,colframe=excol,
  title={\textsc{Corrigé}~\thetcbcounter\ifx\relax#1\relax\relax\else\textmd{ --- \itshape #1}\fi}}
\newcommand{\vv}[1]{\vec{#1}}\newcommand{\ds}{\displaystyle}
\newcommand{\R}{\mathbb{R}}\newcommand{\N}{\mathbb{N}}\newcommand{\Z}{\mathbb{Z}}
% (un \usepackage{fancyhdr}+en-tête est possible ; il est ignoré côté web)
% =========================================================================
\begin{document}
\sisetup{per-mode=symbol}

\begin{corrige}[travail du poids sur un plan incliné]
\begin{enumerate}[label=\alph*)]
  \item $h=AB\times\sin\ang{30}=6\times0,5=\qty{3}{\metre}$.
  \item Poids : $G=mg=20\times10=\qty{200}{\newton}$.
        \begin{itemize}[leftmargin=1.4em,topsep=1pt]
          \item par la hauteur : $W(\vv G)=mgh=20\times10\times3=\qty{600}{\joule}$ ;
          \item par la composante : $W(\vv G)=(G\sin\ang{30})\times AB=(200\times0,5)\times6=\qty{600}{\joule}$.
        \end{itemize}
        Les deux méthodes donnent le même résultat : $W(\vv G)=\qty{600}{\joule}$ (travail moteur, la
        descente accélère l'objet).
\end{enumerate}
\end{corrige}

\begin{corrige}[puissance d'un moteur à vitesse constante]
\begin{enumerate}[label=\alph*)]
  \item $v=\dfrac{72}{3,6}=\qty{20}{\metre\per\second}$.
  \item À vitesse constante, la force motrice est dirigée dans le sens du mouvement et vaut
        $F=F_f=\qty{400}{\newton}$. Donc
        $P=F\times v=400\times20=\qty{8000}{\watt}=\qty{8}{\kilo\watt}$.
\end{enumerate}
\end{corrige}

\begin{corrige}[travail comme aire sous la courbe $F(x)$]
L'aire d'un trapèze vaut $\dfrac{(\text{grande base}+\text{petite base})}{2}\times\text{largeur}$ :
\[
W=\frac{F_A+F_B}{2}\times(x_B-x_A)=\frac{10+30}{2}\times4=20\times4=\qty{80}{\joule}.
\]
Le travail d'une force variable est bien l'aire algébrique sous la courbe $F(x)$.
\end{corrige}

\begin{corrige}[bilan des travaux sur un trajet horizontal]
\begin{enumerate}[label=\alph*)]
  \item $W(\vv F)=F\times d\times\cos\ang{60}=100\times20\times0,5=\qty{1000}{\joule}$.
  \item $W(\vv F_f)=F_f\times d\times\cos\ang{180}=25\times20\times(-1)=\qty{-500}{\joule}$.
  \item Le poids (vertical) et la réaction normale (verticale) sont \emph{perpendiculaires} au
        déplacement horizontal : $W(\vv G)=W(\vv N)=\qty{0}{\joule}$.
  \item $\sum W=1000+(-500)+0+0=\qty{500}{\joule}$ : le travail total reçu est positif, la luge
        accélère.
\end{enumerate}
\end{corrige}

\begin{corrige}[puissance moyenne d'une grue]
\begin{enumerate}[label=\alph*)]
  \item $W=mgh=500\times10\times12=\qty{60000}{\joule}=\qty{60}{\kilo\joule}$.
  \item $P=\dfrac{W}{\Delta t}=\dfrac{60000}{8}=\qty{7500}{\watt}=\qty{7,5}{\kilo\watt}$.
\end{enumerate}
\end{corrige}
\end{document}
